% ========== 2.6 CONSTRAINTS AND ASSUMPTIONS ==========
% Project limitations and assumptions

\section{2.6 Constraints and Assumptions}
\label{sec:constraints-assumptions}

\lettrine{C}{onstraints and assumptions} define the boundaries within which the system must operate and the conditions assumed to be true during development and deployment. Understanding these factors is essential for setting realistic expectations and making appropriate design decisions.

Technical constraints limit the technologies and approaches available for system development. The system must operate on standard developer workstations without requiring specialized hardware, though it may leverage available accelerators when present. Cross-platform compatibility requirements constrain technology choices to those supporting Windows, macOS, and Linux environments. Integration with existing development tools and workflows limits the degree to which the system can deviate from established conventions and standards.

Resource constraints affect the scope and timeline of system development. The project operates within the context of a graduation project with limited development time and team size. This constraint necessitates careful prioritization of features and capabilities, focusing on core functionality that demonstrates the viability of the approach rather than attempting comprehensive implementation of all possible features. External dependencies on AI services and models introduce constraints related to availability, performance, and cost.

\textbf{Key Constraints:}
\begin{compactlist}
    \item Development timeline limited to academic project schedule
    \item Team size constrained to five developers with varying expertise
    \item Hardware requirements limited to standard developer workstations
    \item Dependency on external AI services for model capabilities
\end{compactlist}

Several assumptions underlie the system requirements and design. The project assumes that users have basic programming knowledge and familiarity with development tools, eliminating the need for extensive introductory tutorials. It assumes reliable internet connectivity for accessing external AI services, though local operation should remain possible with reduced functionality. The system assumes that users work primarily with text-based programming languages rather than visual or domain-specific languages.

Assumptions about user behavior and preferences influence requirement priorities. The project assumes that developers value productivity and efficiency over extensive customization options, preferring sensible defaults to complex configuration. It assumes that users are willing to adopt new interaction patterns if they provide clear benefits, but will resist changes that disrupt established workflows without obvious advantages. These assumptions guide decisions about feature implementation and user interface design throughout the project.
